% Recovered verbatim from chapters_16_20_with_part/chapters_16_20_with_part.tex
% Canonical chapter 20
\chapter{Tensor Networks and Holographic Codes}

\section{Tensor Network Realisation}

A tensor network realisation of MK--43 Ultra associates a tensor to each node of the ER graph, with legs corresponding to ER links and physical indices corresponding to local degrees of freedom.
The chromatic fields determine bond dimensions or weights, while the phase field modulates the network structure and boundary conditions.

\section{Holographic Quantum Error-Correcting Codes}

In holographic regimes, the tensor network can be interpreted as a quantum error-correcting code, with bulk logical information encoded redundantly in boundary degrees of freedom.
The ER structure then enforces code constraints, and the white-state core can be seen as a fixed point of an entanglement renormalisation flow.

\section{Towards a Unified Categorical Picture}

The tensor network and code perspectives admit reformulation in higher-categorical language:
objects as sectors (phase, chromatic, ER, geometric), 1-morphisms as couplings, and 2-morphisms as transformations between couplings.
This points toward a unified categorical formulation of MK--43 Ultra, where geometry, information, and dynamics are different facets of a single structured object.
